% ------------------------------------------------------------- LỜI CẢM ƠN
\chapter*{Lời cảm ơn}
\addcontentsline{toc}{chapter}{Lời cảm ơn}
\markboth{Lời cảm ơn}{}

Em xin gửi lời cảm ơn chân thành tới \textbf{TS. Nguyễn Thị Thu Hương} – giảng viên hướng dẫn học phần Project II – đã định hướng đề tài, góp ý về phương pháp phân tích thiết kế hướng đối tượng và cách tổ chức tài liệu trong suốt quá trình thực hiện.

Em cũng xin cảm ơn các thầy cô Trường Công nghệ Thông tin và Truyền thông, Đại học Bách khoa Hà Nội đã trang bị những kiến thức nền tảng về công nghệ phần mềm, cơ sở dữ liệu và lập trình ứng dụng web – là cơ sở để em hoàn thành đề tài này.

Do thời gian và kinh nghiệm còn hạn chế, báo cáo không tránh khỏi thiếu sót. Em rất mong nhận được ý kiến đóng góp của cô để hoàn thiện sản phẩm.

\vspace{0.6cm}
\begin{flushright}
  Hà Nội, tháng 11 năm 2026\\
  Sinh viên thực hiện\\[1.4cm]
  \textbf{Nguyễn Lê Quân}
\end{flushright}

% ------------------------------------------------------------- TÓM TẮT
\chapter*{Tóm tắt nội dung}
\addcontentsline{toc}{chapter}{Tóm tắt nội dung}
\markboth{Tóm tắt nội dung}{}

Đề tài 1A yêu cầu phân tích, thiết kế hệ thống hỗ trợ quản lý một trung tâm luyện thi chứng chỉ ngoại ngữ (IELTS, TOEIC, VSTEP, Cambridge) và lập trình bốn chức năng: \textit{quảng bá thông tin lớp học}, \textit{nhận đăng ký từ học viên}, \textit{tra cứu danh sách lớp} và \textit{chat room tư vấn}. Điểm khó của đề bài nằm ở các ràng buộc nghiệp vụ: học viên chỉ được kiểm tra trình độ sau khi nộp học phí; lớp Online mở khi có từ 10 học viên; lớp Offline mở khi đủ 20, tối đa 30 và khi vượt 30 thì việc tách lớp thuộc quyền Lãnh đạo trung tâm.

Báo cáo trình bày toàn bộ quy trình: (1) bóc tách đề bài thành 14 yêu cầu chức năng và 8 quy tắc nghiệp vụ được đánh mã; (2) mô hình hoá bằng UML với 21 biểu đồ (use case, hoạt động, tuần tự, trạng thái, lớp, ERD, kiến trúc) được viết dưới dạng mã nguồn và render cục bộ; (3) thiết kế cơ sở dữ liệu SQLite 14 bảng chuẩn 3NF, trong đó sĩ số lớp là giá trị \emph{tính toán} và các ràng buộc quan trọng được cưỡng chế ngay ở mức CSDL; (4) cài đặt ứng dụng web Node.js/Express/Socket.IO gồm cổng học viên và dashboard quản trị cho ba vai trò Tư vấn viên, Giáo vụ, Lãnh đạo; (5) kiểm thử tự động 78 ca kiểm thử (đơn vị, tích hợp API, chat realtime) đạt 100\%.

Một kết quả phân tích đáng chú ý: khi lớp Offline vượt 30 học viên, phương án “chia đôi” chỉ hợp lệ nếu tổng sĩ số đạt ít nhất 40, vì chia 31 học viên thành 16 + 15 sẽ tạo ra hai lớp không đủ điều kiện mở. Hệ thống vì vậy tự động đề xuất phương án phù hợp (chia đều, tách phần dư hoặc giữ ngoại lệ theo sức chứa phòng) để Lãnh đạo quyết định.

\noindent\textbf{Từ khoá:} phân tích thiết kế hướng đối tượng, UML, quản lý trung tâm ngoại ngữ, Node.js, Socket.IO, SQLite.
